% sections/5_conclusions.tex
\section{Conclusiones}

El desarrollo computacional de este trabajo ha permitido simular y caracterizar con alta fidelidad la transición de fase ferromagnética-paramagnética del modelo de Ising bidimensional. Mediante el algoritmo de Monte Carlo, se ha verificado la rotura espontánea de simetría en la vecindad de la temperatura crítica de Onsager ($T_c \approx 2.269$). Además, el análisis de escalado de tamaño finito (\textit{finite-size scaling}) ha corroborado las divergencias teóricas con notable precisión, obteniendo un cociente de exponentes críticos $\gamma/\nu \approx 1.814$. 

Por otro lado, la incorporación del campo magnético externo ha modelado con éxito la dinámica fuera del equilibrio termodinámico reversible. La simulación del ciclo de histéresis en el régimen frío ($T=0.25$) reproduce la memoria magnética y la metaestabilidad del material, evidenciando cómo la reinversión estructural opera como una transición de primer orden desencadenada abruptamente por la competencia entre la interacción ferromagnética y el Hamiltoniano de Zeeman.

Finalmente, la integración de técnicas de aprendizaje automático ha demostrado la viabilidad operativa de los modelos de clasificación para trazar fronteras de fases termodinámicas. Aunque la regresión logística aplicada sobre observables macroscópicos precalculados ($M, E, \chi$) reduce el problema a una optimización analíticamente trivial, el clasificador alcanzó una exactitud del 98\%, localizando el umbral crítico de forma robusta. Este ejercicio pedagógico sienta una base computacional sólida para futuras líneas de investigación, justificando la transición hacia arquitecturas de aprendizaje profundo (\textit{Deep Learning}) alimentadas con microestados espaciales crudos, cuyo objetivo será el descubrimiento de nuevas fases en sistemas donde el parámetro de orden sea desconocido \textit{a priori}.